\chapter{آشنایی با داده و پیش‌پردازش}
\label{ch:data}

\section*{اهداف این فصل}
\begin{itemize}
    \item شناخت انواع ویژگی‌ها و مقیاس‌های اندازه‌گیری داده
    \item درک مفهوم کیفیت داده و منابع اصلی مشکلات آن
    \item تسلط بر روش‌های پاک‌سازی، یکپارچه‌سازی، تبدیل و کاهش داده
    \item آشنایی با معیارهای شباهت و فاصله میان نمونه‌ها
\end{itemize}

\section{انواع ویژگی‌ها}
\label{sec:attribute-types}

هر مجموعه داده از تعدادی \dmterm{\textbf{نمونه}}{Instance / Object / Record}
تشکیل شده که هر یک با مجموعه‌ای از \dmterm{\textbf{ویژگی‌ها}}{Attributes / Features}
توصیف می‌شود. شناخت دقیق نوع هر ویژگی، پیش‌نیاز انتخاب روش پیش‌پردازش و الگوریتم
مناسب است. ویژگی‌ها را می‌توان از دو منظر دسته‌بندی کرد.

\subsection{دسته‌بندی کیفی در برابر کمّی}

\begin{itemize}
    \item \dmterm{\textbf{ویژگی‌های کیفی}}{Categorical / Qualitative}
    \begin{itemize}
        \item \dmterm{\textit{اسمی}}{Nominal}: مقادیر فقط برچسب‌اند و ترتیب یا فاصلهٔ
              معناداری ندارند. مثال: رنگ چشم، شهر محل سکونت.
        \item \dmterm{\textit{ترتیبی}}{Ordinal}: مقادیر ترتیب معناداری دارند، اما فاصلهٔ
              میان آن‌ها مشخص یا یکسان نیست. مثال: سطح تحصیلات (دیپلم، کارشناسی،
              کارشناسی‌ارشد)، رضایت مشتری (کم، متوسط، زیاد).
    \end{itemize}
    \item \dmterm{\textbf{ویژگی‌های کمّی}}{Numeric / Quantitative}
    \begin{itemize}
        \item \dmterm{\textit{بازه‌ای}}{Interval}: فاصله میان مقادیر معنادار است، اما
              صفر مطلق ندارد. مثال: دما بر حسب سلسیوس.
        \item \dmterm{\textit{نسبتی}}{Ratio}: علاوه بر فاصلهٔ معنادار، صفر مطلق نیز دارد،
              بنابراین نسبت دو مقدار هم معنادار است. مثال: سن، درآمد، وزن.
    \end{itemize}
\end{itemize}

\subsection{دسته‌بندی گسسته در برابر پیوسته}

از منظری دیگر، ویژگی‌ها به \textbf{گسسته} (تعداد محدود یا شمارش‌پذیر از
مقادیر، مانند تعداد فرزندان) و \textbf{پیوسته} (بی‌نهایت مقدار ممکن در یک
بازه، مانند قد یا وزن) تقسیم می‌شوند. تشخیص درست این نوع، در انتخاب روش
گسسته‌سازی (بخش \ref{subsec:discretization}) و همچنین نوع نمودار مناسب برای
تحلیل اکتشافی داده اهمیت دارد.

\section{کیفیت داده}
\label{sec:data-quality}

پیش از هر کاوشی، باید کیفیت داده ارزیابی شود. مهم‌ترین مشکلات رایج در داده‌های
واقعی عبارت‌اند از:

\begin{itemize}
    \item \dmterm{\textbf{مقادیر گمشده}}{Missing Values}: به دلایلی مانند خطای ثبت،
          عدم پاسخ‌گویی کاربر یا خرابی حسگر.
    \item \dmterm{\textbf{نویز}}{Noise}: خطای تصادفی یا انحراف در مقدار اندازه‌گیری‌شده.
    \item \dmterm{\textbf{داده‌های ناسازگار}}{Inconsistent Data}: تناقض میان رکوردها،
          مثلاً دو مقدار متفاوت برای سن یک فرد در دو جدول مختلف.
    \item \dmterm{\textbf{داده‌های تکراری}}{Duplicate Data}: ثبت چندبارهٔ یک رکورد.
    \item \dmterm{\textbf{داده‌های پرت}}{Outliers}: مقادیری که به‌طور قابل‌توجهی با
          بقیهٔ داده تفاوت دارند (جزئیات بیشتر در فصل \ref{ch:outliers}).
\end{itemize}

نکتهٔ مهم این است که نویز و داده پرت یکسان نیستند: نویز معمولاً خطای تصادفی
و بی‌معنی است که باید حذف یا اصلاح شود، اما داده پرت می‌تواند خودِ هدف تحلیل
باشد (مثلاً تراکنش متقلبانه).

\section{پیش‌پردازش داده}
\label{sec:preprocessing}

\subsection{پاک‌سازی داده}

رایج‌ترین راهکارها برای مقادیر گمشده عبارت‌اند از:

\begin{itemize}
    \item حذف رکوردهایی که مقدار گمشده دارند (فقط زمانی که تعداد آن‌ها کم باشد).
    \item جایگذاری با میانگین یا میانه (برای ویژگی‌های عددی) یا مد (برای
          ویژگی‌های کیفی).
    \item جایگذاری با استفاده از یک مدل پیش‌بینی (مثلاً رگرسیون یا نزدیک‌ترین
          همسایه) بر اساس سایر ویژگی‌ها.
\end{itemize}

برای هموارسازی نویز نیز روش‌هایی مانند \dmterm{\textbf{بین‌بندی}}{Binning}،
رگرسیون و تحلیل خوشه‌ای (برای شناسایی و حذف نقاط دورافتاده) به کار می‌روند.

\subsection{یکپارچه‌سازی داده}

وقتی داده از چند منبع مختلف (مثلاً چند پایگاه داده یا چند فایل) جمع‌آوری
می‌شود، باید به مسائلی مانند موارد زیر توجه کرد:

\begin{itemize}
    \item \dmterm{\textbf{تشخیص موجودیت یکسان}}{Entity Identification}: آیا
          \texttt{customer\_id} در یک جدول همان \texttt{cust\_no} در جدول
          دیگر است؟
    \item \dmterm{\textbf{افزونگی}}{Redundancy}: آیا یک ویژگی از ترکیب چند ویژگی
          دیگر قابل استخراج است؟ (مثلاً سود، از تفاضل درآمد و هزینه.)
    \item \textbf{تناقض مقیاس یا واحد}: مثلاً دما در یک منبع بر حسب سلسیوس و
          در منبع دیگر بر حسب فارنهایت ثبت شده باشد.
\end{itemize}

\subsection{تبدیل و نرمال‌سازی داده}

بسیاری از الگوریتم‌ها (به‌ویژه آن‌هایی که بر پایهٔ فاصله عمل می‌کنند، مانند
نزدیک‌ترین همسایه یا \LR{k-means}) به مقیاس ویژگی‌ها حساس‌اند. اگر یک ویژگی در
بازهٔ $[0, 1]$ و ویژگی دیگر در بازهٔ $[0, 100000]$ باشد، ویژگی دوم به‌طور
نادرست بر محاسبهٔ فاصله غالب می‌شود. دو روش رایج نرمال‌سازی عبارت‌اند از:

\paragraph{نرمال‌سازی \LR{Min-Max}}
مقدار هر ویژگی به بازهٔ $[0, 1]$ نگاشت می‌شود:
\[
x' = \frac{x - \min(x)}{\max(x) - \min(x)}
\]

\paragraph{استانداردسازی \LR{Z-score}}
مقدار هر ویژگی بر اساس میانگین ($\mu$) و انحراف معیار ($\sigma$) آن
تبدیل می‌شود:
\[
x' = \frac{x - \mu}{\sigma}
\]

استانداردسازی \LR{Z-score} معمولاً در حضور داده‌های پرت پایدارتر از
\LR{Min-Max} عمل می‌کند، زیرا \LR{Min-Max} به شدت به مقادیر حداقل و حداکثر
حساس است.

\subsection{کاهش داده و کاهش ابعاد}

هدف کاهش داده، رسیدن به بازنمایی کوچک‌تر اما تقریباً برابر از نظر محتوای
اطلاعاتی است. دو رویکرد اصلی عبارت‌اند از:

\begin{itemize}
    \item \textbf{کاهش تعداد نمونه‌ها}: از طریق \dmterm{نمونه‌برداری}{Sampling} یا
          تجمیع داده.
    \item \textbf{کاهش تعداد ویژگی‌ها (کاهش ابعاد)}: از طریق
          \dmterm{انتخاب زیرمجموعه‌ای از ویژگی‌ها}{Feature Selection} یا
          استخراج ویژگی جدید.
\end{itemize}

رایج‌ترین روش استخراج ویژگی، \dmterm{\textbf{تحلیل مؤلفه‌های اصلی}}{Principal Component Analysis (PCA)}
است. ایدهٔ اصلی \LR{PCA} یافتن جهت‌هایی (مؤلفه‌های اصلی) در فضای ویژگی‌هاست که
بیشترین واریانس داده را حفظ می‌کنند؛ با تصویر کردن داده روی چند مؤلفهٔ اول،
می‌توان ابعاد را کاهش داد در حالی که بخش عمدهٔ اطلاعات داده حفظ می‌شود. توجه
شود که \LR{PCA} پیش از اجرا نیازمند استانداردسازی داده است، زیرا بر پایهٔ
واریانس عمل می‌کند.

\subsection{گسسته‌سازی و سلسله‌مراتب مفهومی}
\label{subsec:discretization}

برخی الگوریتم‌ها (مانند نسخه‌های پایهٔ درخت تصمیم یا بیز ساده) بهتر با
داده‌های کیفی کار می‌کنند. \textbf{گسسته‌سازی} فرآیند تبدیل یک ویژگی عددی
پیوسته به تعدادی بازهٔ گسسته (مانند تبدیل سن به گروه‌های «جوان، میان‌سال،
مسن») است. روش‌های رایج شامل بین‌بندی با عرض یکسان\footnote{\LR{Equal-width}}،
بین‌بندی با فراوانی یکسان\footnote{\LR{Equal-depth}} و روش‌های مبتنی بر
خوشه‌بندی هستند.
\dmterm{\textbf{سلسله‌مراتب مفهومی}}{Concept Hierarchy} نیز به تعریف سطوح مختلف
انتزاع برای یک ویژگی اشاره دارد (مثلاً شهر $\to$ استان $\to$ کشور)، که در
تحلیل‌های چندسطحی کاربرد دارد.

\section{معیارهای شباهت و فاصله}
\label{sec:similarity}

بسیاری از الگوریتم‌های داده‌کاوی (خوشه‌بندی، نزدیک‌ترین همسایه، تشخیص داده پرت)
بر پایهٔ سنجش شباهت یا فاصله میان دو نمونه عمل می‌کنند. رایج‌ترین معیارها
عبارت‌اند از:

\paragraph{فاصلهٔ اقلیدسی}
برای دو بردار $x$ و $y$ در فضای $n$-بعدی:
\[
d(x, y) = \sqrt{\sum_{i=1}^{n} (x_i - y_i)^2}
\]

\paragraph{فاصلهٔ منهتن}
\[
d(x, y) = \sum_{i=1}^{n} |x_i - y_i|
\]

\paragraph{فاصلهٔ مینکوفسکی}
تعمیمی از دو فاصلهٔ بالا با پارامتر $p$:
\[
d(x, y) = \left( \sum_{i=1}^{n} |x_i - y_i|^p \right)^{1/p}
\]
که برای $p=1$ به فاصلهٔ منهتن و برای $p=2$ به فاصلهٔ اقلیدسی تبدیل می‌شود.

\paragraph{شباهت کسینوسی}
برای داده‌های با ابعاد بالا و پراکنده (مانند بردارهای متنی)، شباهت کسینوسی
که زاویهٔ میان دو بردار را می‌سنجد، اغلب مناسب‌تر از فاصلهٔ اقلیدسی است:
\[
\text{sim}(x, y) = \frac{x \cdot y}{\|x\| \, \|y\|}
\]

برای ویژگی‌های اسمی یا دودویی نیز معیارهایی مانند
\dmterm{\textbf{ضریب تطابق ساده}}{Simple Matching Coefficient} و
\dmterm{\textbf{ضریب ژاکارد}}{Jaccard Coefficient} به کار می‌روند که در فصل
مربوط به خوشه‌بندی (فصل \ref{ch:clustering}) دوباره به آن‌ها پرداخته می‌شود.

\section*{پیاده‌سازی در پایتون}

نمونه‌کد کامل این فصل — شامل جایگذاری مقادیر گمشده، استانداردسازی با
\texttt{StandardScaler}، کاهش ابعاد با \texttt{PCA} و محاسبهٔ فاصلهٔ
اقلیدسی/کسینوسی با \texttt{scipy} — در بخش «فصل ۲» پیوست الف
(\ref{app:python}) آمده است.

\section*{خلاصهٔ فصل}

در این فصل با انواع ویژگی‌ها، منابع اصلی مشکلات کیفیت داده، و مراحل کلیدی
پیش‌پردازش شامل پاک‌سازی، یکپارچه‌سازی، تبدیل، کاهش داده و گسسته‌سازی آشنا
شدیم. همچنین رایج‌ترین معیارهای فاصله و شباهت را که در فصل‌های بعدی
(به‌ویژه طبقه‌بندی و خوشه‌بندی) بارها استفاده خواهند شد، مرور کردیم. در
فصل بعد، به اولین وظیفهٔ اصلی داده‌کاوی یعنی طبقه‌بندی می‌پردازیم.

\section*{تمرین‌ها}

\begin{enumerate}
    \item برای هر یک از ویژگی‌های «شمارهٔ ملی»، «دمای بدن»، «رتبهٔ مسابقه»
          و «تعداد فرزندان» نوع آن (اسمی/ترتیبی/بازه‌ای/نسبتی) را مشخص کنید.
    \item فرض کنید در یک مجموعه داده، ۱۵ درصد مقادیر ستون «درآمد» گمشده است.
          دو راهکار برای مواجهه با این مشکل پیشنهاد دهید و مزایا و معایب
          هرکدام را توضیح دهید.
    \item تفاوت نرمال‌سازی \LR{Min-Max} و استانداردسازی \LR{Z-score} را با
          یک مثال عددی نشان دهید و توضیح دهید کدام‌یک در حضور داده پرت
          پایدارتر است.
    \item فاصلهٔ اقلیدسی و فاصلهٔ منهتن را برای دو بردار
          $x = (2, 4, 6)$ و $y = (5, 1, 8)$ محاسبه کنید.
    \item چرا شباهت کسینوسی برای مقایسهٔ اسناد متنی معمولاً بهتر از فاصلهٔ
          اقلیدسی عمل می‌کند؟
\end{enumerate}
